\section{Taxonomy of Reported Evidence}

The evidence is organized along seven questions rather than a single ranking: (1) what compositions and structures have been reported, (2) how powders or crystals were produced, (3) how phase purity was established, (4) which kinetic pathways retain metastable products, (5) what near-neighbor systems can transfer, (6) which thermodynamic data constrain the composition space, and (7) how these observations translate into an experimental design. This faceted organization follows the distinction between composition, process, and measurement emphasized in the supplied taxonomy \cite{chen2024navigating,wang2024optimal,arroyave2022phase}.

\subsection{Target compositions and structures}

Four-component Ba--Y--Zn--Si--O records are limited to glasses or precursor formulations, whereas Ba--Y, Ba--Zn, and Y--Zn silicates provide incomplete-component structures \cite{kozhevnikova2024synthesis,kaur2011influence,sayyed2026physical,kolitsch2006bay2si3o10,kaiser2002crystal,bordun2011cathodoluminescence}. The boundary is compositional: a structurally similar silicate is a neighbor, not a target-phase observation.

\subsection{Synthesis and crystal-growth routes}

Solid-state and combustion routes expose precursor reactivity and calcination effects \cite{shen2019synthesis,faat2023enhanced,richhariya2021synthesis,takesue2009thermal}. Flux growth and solution-mediated routes provide access to single crystals or fine-scale mixing, but flux chemistry and volatility must be retained as experimental variables \cite{bugaris2012materials,goncalves2025flux,kumar1974growth,smith1974flux,ranganathan2008sol,parsa2018synthesis}. Thermal-profile studies form a separate branch because crucible, atmosphere, holding time, and cooling rate can change the product even at fixed nominal composition \cite{thieme2022noble,xu2020crystallization,shaaban2021thermal,zheng2025effect}.

\subsection{Characterization and phase purity}

Powder or single-crystal diffraction establishes long-range order, but minor phases and composition gradients may be invisible in a single pattern. Quench--XRD, electron-probe mapping, and in-situ microscopy provide complementary views of phase assemblage and nucleation \cite{kim2025key,derensy2026phase,martinson2019pellet,abdeyazdan2024phase,mckenzie2021nucleation,bocker2012in}. Spectroscopy and local-structure probes are especially useful when glass-ceramic products contain overlapping diffraction features \cite{mastelaro2018x,abdelhameed2023structural,srabionyan2024local,shaaban2021spectroscopic}.

\subsection{Kinetics and thermodynamics}

Glass nucleation, mechanochemical intermediates, interfacial melts, and finite holds constitute kinetic branches that cannot be read directly from an equilibrium diagram \cite{moulton2021a,rana2021insights,schmelzer2020effects,zhang2026interfacial,jang2025identification}. Binary and ternary phase-equilibrium studies provide quantitative thermodynamic anchors, while multicomponent databases and autonomous-materials methods provide strategies for navigating sparse data \cite{shevchenko2021integrated,jeon2025phase,kim2025key,chen2024navigating,szymanski2023an}.

The taxonomy therefore encodes transfer boundaries explicitly: a paper can support a structural motif, a route concept, or a thermodynamic constraint, but none of these alone establishes a phase-pure BYZSO quaternary. This rule prevents near-neighbor evidence from being mistaken for direct target evidence.
